In this work, we modeled memory consolidation as slow, reactivation-driven dynamics that reorganizes brain-wide engram patterns over slow timescales of hours to years. By generalizing the classical area-centered view of systems consolidation to a global population-level description, we revealed how established recurrent network motifs can be reinterpreted as mechanisms of long-term memory reorganization.
Representational drift emerges as a natural corollary of this global memory reorganization. Finally, subsampling---the observation of a small fraction of the full population---can obscure directed consolidation dynamics, such that they appear random, complicating the recovery of the system dynamics from neuronal recordings. 

The suggested dynamical systems perspective on memory consolidation creates a bridge from classical theories of consolidation, which primarily describe the time-course of hippocampal dependency and differ in their postulates for episodic and semantic memories \citep{Squire1995,Nadel1997,Winocur2010}, to contemporary views of global memory organization \citep{Moscovitch2022}. The area-centered view is captured as a special case of a general brain-wide description of memory reorganization. A sustained or fading hippocampal dependency is described by the geometric alignment of memories with the consolidation process, represented by the consolidation matrix~$A$: Depending on how an engram aligns with this matrix, it is passively forgotten, actively maintained within the hippocampus, or consolidated into the neocortex. The modeling framework does not provide a natural distinction between episodic and semantic memory and therefore cannot speak to their differential consolidation. However, it can readily describe a selective and partial consolidation such that memories could be stripped of their episodic vividness, rendering them increasingly semantic in nature. 

\subsection{Mechanistic interpretation \& assumptions}

The suggested model abstracts away mechanistic details at the cellular and circuit levels and instead focuses on the phenomenological evolution of memory traces. It distinguishes between active consolidation and passive decay, treating both as deterministic drivers of engram evolution. 
The decay term represents the aggregate effect of stochastic, destructive processes, such as interference from newly stored memories \citep{Robertson2012}, spontaneous synaptic turnover \citep{Ziv2018}, and neurogenesis \citep{Golbabaei2025a}. We note that this formulation captures physical trace erosion rather than forgetting due to transient retrieval failure \citep{deSnoo2025}.

The active consolidation component represents the systematic drive from processes such as circuit-level engram transfer \citep{McClelland1995,Remme2021,Bhasin2024}, synaptic plasticity \citep{Lee2023}, assembly formation \citep{Tome2024,Golbabaei2025}, and synaptic consolidation \citep{Fusi2005,Benna2016}, all of which dictate how a memory trace reorganizes over time. The effect of these mechanisms is phenomenologically summarized in the consolidation matrix~$\mathbf{A}$, so we assumed fixed, linear engram dynamics. This linearity limits the model, because all memories undergo the same dynamics and different aspects of a memory cannot interact nonlinearly. Biological consolidation, however, is adaptive and nonlinear. Experience can reshape the consolidation dynamics themselves in a process known as schema formation \citep{Tse2007} or meta-learning \citep{Wang2018,Zhou2025}. Similarly, the consolidation dynamics often depend on the consolidation context, such as internal states \citep{Chouhan2020}, or the systematic interaction with other memories \citep{McClelland1995,Schlichting2015}. Our linear framework thus represents a first-order approximation: it captures how memories evolve in a given consolidation context when the underlying consolidation rules are held constant. Generalization to nonlinear dynamics is beyond the scope of the present work. 

The assumption that the consolidation of an engram depends on the engram itself amounts to assuming that consolidation is driven by reactivation, because information about the shape of the engram is assumed to be present at the time of consolidation. We remain agnostic as to the type of reactivation involved. In particular, we do not distinguish between replay during sleep \citep{Diekelmann2010}, wakeful rest \citep{Wamsley2022}, and the re-experience of a known learning paradigm \citep{Himmer2019}, even though they might contribute differentially to consolidation. Finally, the model treats reactivation merely as a driver of memory reorganization and ignores a potential re-encoding of memory traces during re-experience \citep{Nadel1997}, systems reconsolidation \citep{Debiec2002}, or direct interactions between distinct memory traces beyond passive interference \citep{Schlichting2015}.

\subsection{Representational drift as a correlate of consolidation}

In our model, the same consolidation dynamics that reorganize memory engrams across the brain over years also produce representational drift in local population codes over days. In this view, representational drift is not merely a passive degradation of neural codes that needs compensation \citep{Rule2020,KalleKossio2021} but a signature of distributed systems consolidation. The model hence unifies the two phenomena by treating the stimulus-behavior mappings in drift studies as cue-response associations that undergo consolidation.

In our model, systems consolidation is deterministic and directed, implying that the associated drift is also directed. Existing models for drift tend to drive it with a stochastic component, attributed to random synaptic updates \citep{Qin2023,Ratzon2024,Morales2024,Eppler2026}, interference from ongoing learning \citep{Devalle2025}, or fluctuations in intrinsic excitability \citep{Delamare2024,Haimerl2025}. In our framework, on the other hand, these same stochastic mechanisms are captured as a passive decay of the engram and do not drive the drift dynamics. Most existing models include an additional directed component in the drift dynamics, e.g.~by constraining the random dynamics to a solution manifold \citep{Ratzon2024,Natrajan2024}, balancing it with plasticity \citep{Qin2023,Devalle2025,Morales2024,Eppler2026}, or by deriving from it a directed area-to-area consolidation component \citep{Kossio2025}. Our model illustrates that a directed component alone---systems consolidation---is sufficient to generate the complex phenomenology of representational drift.

The apparent randomness of drift observed in experiments can be reconciled with deterministic dynamics by accounting for limited observability. When a high-dimensional, deterministic consolidation process is observed through the lens of a small subset of recorded neurons, the resulting dynamics appear dominated by noise and lose the structure and predictability present in the full population. This %illustrates the inherent challenge of inferring global dynamics from local observations \citep{Levina2022} and 
suggests that random components of drift reported in experimental recordings \citep{Rule2020,Schoonover2021} may be the result of limited observability rather than a fundamental property of the neural code.

The model predicts that drift dynamics are deterministic and hence predictable. Unfortunately, the experimental test of this prediction requires overcoming the subsampling barrier \citep{Wilting2018,Levina2022,Qian2024} to characterize the full consolidation dynamics. For large neural systems such as mammalian brains, this is challenging because longitudinal recordings only observe a fraction of the full brain over a limited time window. Given such limited data, a linear fit of the representational dynamics is prone to overfitting, and even a failure to fit would be a poor falsification of the model, because it may merely be the result of subsampling. However, we expect that fitting may be possible even in large memory systems if the dimension of the neural subspace covered by the relevant engrams over time is small or comparable to the number of recorded neurons. In this case, the unobserved part of the engram does not contain additional information, and the fitting should succeed. 

 A test of the model could also be within reach for smaller model systems. For example, in the \textit{Drosophila} mushroom body, a small number of mushroom body output neurons represent odor-reward associations \citep{Owald2015b}, and their involvement changes through consolidation \citep{Cervantes-Sandoval2013}. While recording the full population of these neurons simultaneously is not currently possible in Drosophila, this or other model systems may, in the future, offer the required access to study whether memory consolidation, engram dynamics and representational drift are different sides of the same coin. 
